\section*{Capítulo \ref{ch:b2:climate-variability} --- Variabilidade climática}

\begin{solution}{exo:b2:climate-variability:1}
(a) $29.27 - 26.71 = \qty{2.56}{\celsius}$ em 2026; $28.79 - 26.73 = \qty{2.06}{\celsius}$ em 1997.
(b) Setembro de 2026 foi ao mesmo tempo o mês mais quente, por \qty{0.48}{\celsius}, e a maior anomalia, por
\qty{0.50}{\celsius}. (c) Cada ano é comparado com um período de base de 30 anos centrado nele, deslocado
a cada cinco anos; a \omterm{def:b2:climate-variability:modes}{climatologia} de setembro dos anos 1980 e 1990 difere ligeiramente da dos
anos 2010 e 2020.
\end{solution}

\begin{solution}{exo:b2:climate-variability:2}
(a) JJA: $(0.79 + 0.95 + 1.28)/3 = 1.007$; JAS: $(0.95 + 1.28 + 1.55)/3 = 1.260$; ASO: $(1.28 +
1.55 + 1.68)/3 = 1.503$; SON: $(1.55 + 1.68 + 2.01)/3 = 1.747$; OND: $(1.68 + 2.01 + 2.02)/3 =
1.903$, tudo em \unit{\celsius}. (b) Arredondados a uma casa decimal: 1.0, 1.3, 1.5, 1.7, 1.9, os valores da
tabela. (c) Em ASO de 2023, com 1.5.
\end{solution}

\begin{solution}{exo:b2:climate-variability:3}
(a) O índice acompanha o desvio da pressão no Taiti menos a de Darwin: abaixo da média
no Taiti e acima em Darwin, ele é negativo. (b) El~Niño. (c) Alísios de leste mais fracos, até
ventos de oeste perto da superfície no oeste, e menos chuva que o habitual sobre a Indonésia, pois a
convecção profunda deslocou-se para leste.
\end{solution}

\begin{solution}{exo:b2:climate-variability:4}
(a) $804.1/880.6 = \qty{91.3}{\percent}$ da média: entre 90 e \qty{110}{\percent}, uma
estação normal, embora em sua parte inferior. (b) A estação totaliza $556 + 659 + 427 + 252 =
\qty{1894}{mm}$; julho traz $659/1894 = \qty{34.8}{\percent}$ dela, mais de um terço.
\end{solution}

\begin{solution}{exo:b2:climate-variability:5}
(a) $\tau = 1025 \times 3992 \times 50/16.1 = \qty{1.271e7}{s} = \qty{147.1}{\day}$ para
\qty{50}{m}, e o dobro, \qty{294.2}{\day}, para \qty{100}{m}. (b) $\mathrm{e}^{-30.44/147.1} =
0.813$ e $\mathrm{e}^{-30.44/294.2} = 0.902$. (c) O tempo de amortecimento é proporcional a $h$: uma
camada de mistura mais rasa armazena menos calor para o mesmo fluxo e esquece uma anomalia mais depressa, de modo que a
memória do oceano, e a previsibilidade que ela dá, diminuem.
\end{solution}

\begin{solution}{exo:b2:climate-variability:6}
(a) Três meses são 91.3 dias: $1.2 \times \mathrm{e}^{-91.3/147.1} = \qty{0.65}{K}$; seis meses
são 182.6 dias: $1.2 \times \mathrm{e}^{-182.6/147.1} = \qty{0.35}{K}$. (b) A anomalia decai
lentamente, por isso é um palpite melhor que zero para a anomalia do mês seguinte. Mas o forçamento aleatório acrescenta
novas anomalias todos os dias, desconhecidas no momento da previsão; à medida que a anomalia inicial decai,
elas ocupam uma parte cada vez maior, e a previsão deriva para a \omterm{def:b2:climate-variability:modes}{climatologia}, onde já não
ajuda.
\end{solution}

\begin{solution}{exo:b2:climate-variability:7}
(a) $160 \times 86\,400 = \qty{1.38e7}{J}$, \qty{13.8}{MJ} por metro quadrado e por dia. (b) A
camada contém $(1000 - 850) \times 100/9.81 = \qty{1529}{kg}$ de ar por metro quadrado, com uma capacidade
térmica de $1529 \times 1004.5 = \qty{1.536e6}{J.m^{-2}.K^{-1}}$; ela aqueceria $1.38 \times
10^7/1.536 \times 10^6 = \qty{9.0}{K}$ por dia. (c) A camada irradia calor no infravermelho, para o
espaço e para o ar acima; e o ar quente sobe e mistura-se com a camada acima dele, ou é
levado pelos ventos, entre eles a entrada de ar mais fresco que a própria baixa térmica atrai.
\end{solution}

\begin{solution}{exo:b2:climate-variability:8}
(a) $(-1.93 - 1.11 - 1.98)/3 = -1.673$, cerca de $-1.67$. (b) É o inverno mais baixo do registro.
(c) Uma NAO fortemente negativa: uma Baixa da Islândia e uma Alta dos Açores fracas, um jato mais fraco e uma trajetória das tempestades
mais ao sul que o habitual; um inverno frio e seco no norte da Europa, úmido no sul da Europa,
e mais irrupções de ar frio no leste dos Estados Unidos.
\end{solution}

\begin{solution}{exo:b2:climate-variability:9}
(a) $9.3 - 0.4 = \qty{8.9}{mm}$ por dia, mais de 2.5: no domínio. (b) No Hemisfério Sul
o verão é dezembro--fevereiro: $6.1 - 4.2 = \qty{1.9}{mm}$ por dia, menos de 2.5: fora dele,
embora chova ali mais do que na maioria dos lugares do domínio. (c) $1894/122 = \qty{15.5}{mm}$ por
dia; Kerala só ficaria fora se sua taxa de inverno excedesse $15.5 - 2.5 = \qty{13.0}{mm}$ por dia.
(d) Uma \omterm{def:b2:climate-variability:monsoon}{monção} é uma inversão sazonal: um lugar úmido o ano inteiro, como a segunda
estação, não tem \omterm{def:b2:climate-variability:monsoon}{monção}, e um lugar seco com uma estação chuvosa curta e intensa tem uma.
\end{solution}

\begin{solution}{exo:b2:climate-variability:10}
(a) Derivando, $\ddot T = a\dot H = -abT$: a equação de um oscilador harmônico de frequência
angular $\omega = \sqrt{ab}$, e de período $2\pi/\sqrt{ab}$. (b) $2\pi/\sqrt{ab} = 4$ dá $ab =
(2\pi/4)^2 = \qty{2.47}{yr^{-2}}$. (c) Com $T = u\,\mathrm{e}^{-rt}$, $H = v\,\mathrm{e}^{-rt}$,
$\dot T = (\dot u - ru)\,\mathrm{e}^{-rt} = (av - ru)\,\mathrm{e}^{-rt}$, logo $\dot u = av$, e
da mesma forma $\dot v = -bu$: $(u, v)$ oscila sem amortecimento, e $(T, H)$ é essa oscilação
multiplicada por $\mathrm{e}^{-rt}$. Num período a amplitude cai pelo fator $\mathrm{e}^{-0.25 \times
4} = \mathrm{e}^{-1} = 0.37$. (d) Um forçamento: as flutuações aleatórias dos alísios, que
impulsionam o sistema uma e outra vez, e a \omterm{def:b2:climate-variability:bjerknes}{retroalimentação de Bjerknes}, que torna o amortecimento mais fraco ou
até negativo.
\end{solution}

\begin{solution}{exo:b2:climate-variability:11}
(a) $A = R^2 \,\Delta\lambda\,(\sin\ang{5} - \sin(\ang{-5})) = (6.371 \times 10^6)^2 \times 0.8727
\times 0.1743 = \qty{6.17e12}{m^2}$. (b) $Q = 1025 \times 3992 \times 100 \times 2 \times 6.17
\times 10^{12} = \qty{5.05e21}{J}$. (c) Seis meses são $182.6 \times 86\,400 = \qty{1.578e7}{s}$;
por metro quadrado, $8.18 \times 10^8/1.578 \times 10^7 = \qty{51.9}{W/m^2}$, mantidos durante meio
ano. (d) Do próprio oceano: alísios mais fracos fazem subir menos água fria no leste, e a água quente
do oeste, armazenada acima da termoclina mais profunda, espalha-se para leste; grande parte do aquecimento
pode vir de calor que já está no oceano, deslocado e não ganho.
\end{solution}

\begin{solution}{exo:b2:climate-variability:12}
(a) O circuito devolve $0.3g$, depois $0.3g^2$, e assim por diante: o total é $0.3(1 + g + g^2 + \dots) =
0.3/(1 - g)$, uma série geométrica que converge para $g < 1$: \qty{0.6}{K}, \qty{1.5}{K} e
\qty{3.0}{K} para $g = 0.5$, 0.8 e 0.9. (b) A série diverge: a anomalia cresceria sem
limite; o crescimento tem de ser detido por processos que o circuito deixa de fora, o amortecimento pelo fluxo de calor
à superfície e o esvaziamento do estoque de água quente que o oscilador de recarga descreve. (c)
Quanto mais quente a superfície do mar, mais facilmente um pequeno aquecimento desencadeia convecção profunda e chuva,
cujo aquecimento move os ventos; é a estação em que os efeitos do ENSO são mais fortes.
\end{solution}

\begin{solution}{pb:b2:climate-variability:1}
\textbf{1.} $(0.26 + 0.75 + 1.13)/3 = \qty{0.713}{\celsius}$.
\textbf{2.} MJJ $(0.75 + 1.13 + 1.52)/3 = 1.133$; JJA $(1.13 + 1.52 + 1.80)/3 = 1.483$; JAS $(1.52 +
1.80 + 2.06)/3 = 1.793$.
\textbf{3.} ASO $(1.80 + 2.06 + 2.26)/3 = 2.040$; SON $(2.06 + 2.26 + 2.42)/3 = 2.247$; OND $(2.26
+ 2.42 + 2.36)/3 = 2.347$; NDJ $(2.42 + 2.36 + 2.33)/3 = 2.370$.
\textbf{4.} Arredondados a uma casa decimal, 0.7, 1.1, 1.5, 1.8, 2.0, 2.2, 2.3 e 2.4: a tabela.
\textbf{5.} Um único mês traz o ruído do tempo; o ENSO dura muitos meses, e a média de três
elimina boa parte do ruído e conserva o sinal lento.
\textbf{6.} Em JJA de 1997, com 1.483, arredondado para 1.5; atingiu o pico de 2.4 em NDJ.
\textbf{7.} De ASO de 2018 a MJJ de 2019, e de SON de 2019 a FMA de 2020.
\textbf{8.} Dez estações, e seis.
\textbf{9.} Em ASO de 2020; dentro de 2020 só quatro estações, de ASO a NDJ, estão em $-0.5$ ou abaixo, uma
a menos do que exige a regra, e DJF de 2021 a completa.
\textbf{10.} Um El~Niño, de SON de 2018 a MAM de 2019, sete estações (AMJ de 2019 cai para 0.37); a
La~Niña começa em AMJ de 2020 ($-0.55$), quatro estações antes de onde o ONI a situa, em ASO.
\textbf{11.} Não: o índice relativo nunca passa de 0.29 no fim de 2019. Todo o oceano tropical estava
então cerca de meio grau mais quente que seu período de base; a região Niño-3.4 estava quente em
termos absolutos, mas quase nada mais quente que os trópicos à sua volta.
\textbf{12.} Onze.
\textbf{13.} 8, 7, 10, 4, 5, 6, 12, 5, 9 e 3 anos.
\textbf{14.} O mais curto é de 3 anos (2023 a 2026), o mais longo de 12 (1997 a 2009).
\textbf{15.} $69/10 = 6.9$ anos das duas maneiras: a soma dos intervalos sucessivos é o tempo do
primeiro início ao último, $2026 - 1957 = 69$ anos.
\textbf{16.} Ele já durou cinco estações, de MAM a JAS, em \qty{0.5}{\celsius} ou acima, e
atingiu 2.2, acima de 1.5; seu início, MAM de 2026, é a primeira estação da sequência, e estações posteriores
só podem prolongar a sequência, não deslocar seu começo.
\textbf{17.} Não, e não precisa ser: os 25 episódios de todas as intensidades ocorrem a cada três anos, mais ou menos, como dizem as perguntas frequentes;
os fortes são um subconjunto, portanto mais raros, e seus intervalos vão de 3 a 12 anos.
\textbf{18.} $\tau = 1025 \times 3992 \times 50/16.1 = \qty{1.27e7}{s} = \qty{147}{\day}$.
\textbf{19.} Camada de mistura: $1025 \times 3992 \times 50 = \qty{2.05e8}{J.m^{-2}.K^{-1}}$; coluna de ar:
$(10^5/9.81) \times 1004.5 = 10\,194 \times 1004.5 = \qty{1.02e7}{J.m^{-2}.K^{-1}}$, 20 vezes
menos.
\textbf{20.} $1.024 \times 10^7/(1025 \times 3992) = \qty{2.5}{m}$ de água do mar.
\textbf{21.} $1.024 \times 10^7/16.1 = \qty{6.36e5}{s} = 7.4$~dias.
\textbf{22.} Uma anomalia só do ar é amortecida em cerca de uma semana, uma da camada de mistura em
cerca de cinco meses, e a camada quente acima da termoclina é ainda mais profunda: só o oceano guarda
calor por tempo suficiente para levar uma anomalia de uma estação para a seguinte, enquanto a atmosfera fornece
os ventos que acoplam os dois.
\textbf{23.} Desde 1950, os episódios fortes de El~Niño começaram em média a cada \textbf{6.9 anos},
cerca de sete.
\end{solution}
